\section{总结与延伸}

\subsection{核心结论}

本节先给出全篇结论，再保留开放问题。这样读者可以把本期从“资本市场访谈”压缩成一套判断 AI 周期的工作框架。

最后把本期压缩成十条判断。

\begin{enumerate}
\item 2025--2026 的美股科技主线是 AI、再工业化和金融数字化三条线的合流。
\item OpenAI 的关键非共识，是它被看成产品公司，而不只是模型公司。
\item 模型公司现金流短期像负向滚雪球，拐点来自训练成本增速下降或模型 ROI 上升。
\item OpenAI 收入不能只看个人订阅，还要看企业端、API、Agent、广告和电商。
\item Anthropic 和 OpenAI 的分化来自客户结构、产品路线和财务预测假设。
\item Robinhood 展示了周期业务如何通过多元化、份额、定价权和成本控制走出 Alpha。
\item 硅谷资本偏好的“坏小孩”不是任性，而是 founder-led 的反常识执行力。
\item Coding、视频、Agentic Commerce 和垂直应用是 2025 AI 应用最值得研究的赛道。
\item AI 收入早期抢电子收入池，长期要进入劳动力池和新增生产力。
\item 判断 AI bubble 要看用户、收入、ROIC、模型进步和估值，而不是只看市场情绪。
\end{enumerate}

\subsection{开放问题}

最后保留开放问题，是因为本期讨论的许多变量都还在演化：模型收入、Agent Commerce、金融入口、云利润率和 AI 提效财务化，都需要后续季度继续验证。

\begin{itemize}
\item OpenAI 的企业入口会不会比消费端超级入口更早兑现？
\item Anthropic 的 2B/Coding 路线能否保持增速，并形成比 OpenAI 更好的单位经济？
\item Agentic Commerce 会先伤害大型平台，还是先帮助长尾商家？
\item Robinhood 的一站式金融应用能否进入财富管理和另类资产主流市场？
\item AI 应用公司的估值口径会回到毛利率，还是转向单用户毛利和绝对利润？
\item 2026 年 AI 收入是否会成为比 Nvidia 出货更重要的市场锚？
\end{itemize}

\subsection{拓展阅读}

\begin{itemize}
\item EP127 广密跨年大模型季报：AI War、Online Learning 和收入栈。
\item EP136 广密大模型季报第 9 集：Coding、模型 OS 和硅谷御三家。
\item EP139 Agent 综述：Agent 技术史、OpenClaw Moment 与社会辐射。
\item Altimeter/BG2 相关公开访谈：Sam Altman、Jensen Huang 与美国科技资本叙事。
\item OpenAI、Anthropic、Robinhood、Waymo、Polymarket/Kalshi 等公开材料：用于交叉验证收入、产品和监管变化。
\end{itemize}

\begin{importantbox}{最后的判断}
EP125 最值得保留的不是某个估值数字，而是一套资本市场的读法：先问收入池，再问成本曲线；先区分 Beta 和 Alpha，再判断公司执行；先看 AI 投入是否带来收入和 ROIC，再讨论泡沫。这样读 AI 周期，才不容易被单一叙事带跑。
\end{importantbox}

\end{document}
